\documentclass{article}
\usepackage{amsmath}
\usepackage{amssymb}

\title{CSE400 – Fundamentals of Probability in Computing}
\author{Lecture 18: Marginal PMF, Correlation, Covariance, Joint Gaussian Random Variables\\
Instructor: Dhaval Patel, PhD}
\date{March 12, 2026}

\begin{document}

\maketitle

\section{Marginal PMF}

\subsection{Definition}

For discrete random variables $X$ and $Y$, the joint PMF is

\[
p_{XY}(x,y) = P(X=x, Y=y)
\]

The marginal PMFs are obtained by summing over the other variable:

\[
p_X(x) = \sum_y p_{XY}(x,y)
\]

\[
p_Y(y) = \sum_x p_{XY}(x,y)
\]

\textbf{Key Idea}

The marginal PMF of one variable is obtained by summing the joint PMF over all possible values of the other variable.

\subsection{Marginal PMF from a Joint PMF Table}

Row sums give $p_X(x)$ and column sums give $p_Y(y)$.

Let the joint PMF table be:

\[
\begin{array}{c|ccc|c}
p_{XY}(x,y) & y_1 & y_2 & y_3 & p_X(x) \\
\hline
x_1 & p_{11} & p_{12} & p_{13} & \sum_j p_{1j} \\
x_2 & p_{21} & p_{22} & p_{23} & \sum_j p_{2j} \\
x_3 & p_{31} & p_{32} & p_{33} & \sum_j p_{3j} \\
\hline
p_Y(y) & \sum_i p_{i1} & \sum_i p_{i2} & \sum_i p_{i3} &
\end{array}
\]

Probability conditions:

\[
\sum_i p_X(x_i) = 1
\]

\[
\sum_j p_Y(y_j) = 1
\]

\subsection{Example}

Let the joint PMF be

\[
\begin{array}{c|cc}
p_{XY}(x,y) & y=0 & y=1 \\
\hline
x=0 & \frac{1}{8} & \frac{1}{4} \\
x=1 & \frac{3}{8} & \frac{1}{4}
\end{array}
\]

Step 1: Compute $p_X(x)$

\[
p_X(0)=\frac{1}{8}+\frac{1}{4}=\frac{3}{8}
\]

\[
p_X(1)=\frac{3}{8}+\frac{1}{4}=\frac{5}{8}
\]

Step 2: Compute $p_Y(y)$

\[
p_Y(0)=\frac{1}{8}+\frac{3}{8}=\frac{1}{2}
\]

\[
p_Y(1)=\frac{1}{4}+\frac{1}{4}=\frac{1}{2}
\]

Hence

\[
p_X(x)=
\begin{cases}
\frac{3}{8}, & x=0 \\
\frac{5}{8}, & x=1
\end{cases}
\]

\[
p_Y(y)=
\begin{cases}
\frac{1}{2}, & y=0 \\
\frac{1}{2}, & y=1
\end{cases}
\]

\section{Correlation}

\textbf{Definition}

A basic way to measure joint behavior is through the product moment.

In the continuous case,

\[
R_{XY} = E[XY]
\]

It tells us the average value of the product $XY$.

This is often called a raw correlation measure.

\textbf{Interpretation}

This measure:

\begin{itemize}
\item depends on the scale of $X$ and $Y$
\item is computed about the origin
\item mixes together:
\begin{itemize}
\item mean values
\item spread
\item dependence
\end{itemize}
\end{itemize}

Thus it is useful as a starting point but not a clean measure of centered dependence.

If

\[
R_{XY}=0
\]

then the variables are orthogonal about the origin.

\textbf{Transition}

To remove the effect of the means, we center the variables first, which leads to covariance.

\section{Covariance}

\subsection{Definition}

Covariance is defined as

\[
Cov(X,Y)=E[(X-\mu_X)(Y-\mu_Y)]
\]

where

\[
\mu_X = E[X], \quad \mu_Y = E[Y]
\]

\textbf{Interpretation}

We first subtract the mean from each variable.

Then we measure how the centered quantities vary together.

Covariance captures whether deviations from the means occur:

\begin{itemize}
\item in the same direction
\item in opposite directions
\end{itemize}

Interpretation of sign:

\begin{itemize}
\item Positive covariance
\item Negative covariance
\item Zero covariance
\end{itemize}

Covariance tells the direction of joint linear variation between $X$ and $Y$.

However, the magnitude depends on the units of the variables, therefore it cannot be used to assess the strength of the relationship.

\subsection{Algebraic Form and Properties}

Start with

\[
Cov(X,Y)=E[(X-\mu_X)(Y-\mu_Y)]
\]

Expand:

\[
=E[XY-X\mu_Y-\mu_XY+\mu_X\mu_Y]
\]

Apply linearity of expectation:

\[
=E[XY]-\mu_YE[X]-\mu_XE[Y]+\mu_X\mu_Y
\]

Since

\[
\mu_X=E[X], \quad \mu_Y=E[Y]
\]

Therefore

\[
Cov(X,Y)=E[XY]-\mu_X\mu_Y
\]

\subsection{Properties of Covariance}

\begin{enumerate}
\item Symmetry
\[
Cov(X,Y)=Cov(Y,X)
\]

\item Variance as a Special Case
\[
Cov(X,X)=Var(X)
\]

\item Linearity with Respect to Constants
\[
Cov(aX+b,cY+d)=ac\,Cov(X,Y)
\]

\item Independence Implies Zero Covariance
\end{enumerate}

If

\[
X \perp Y
\]

then

\[
Cov(X,Y)=0
\]

Important statement:

Zero covariance means uncorrelated, but does not imply independence in general.

\subsection{Example 1}

If

\[
Y=-2X+3
\]

Find $Cov(X,Y)$.

Using the covariance formula:

\[
Cov(X,Y)=E[XY]-E[X]E[Y]
\]

Since

\[
XY=X(-2X+3)=-2X^2+3X
\]

and

\[
E[Y]=E[-2X+3]=-2E[X]+3
\]

Substitute into covariance formula:

\[
Cov(X,Y)=E[-2X^2+3X]-E[X](-2E[X]+3)
\]

\[
=-2E[X^2]+3E[X]+2(E[X])^2-3E[X]
\]

\[
=-2E[X^2]+2(E[X])^2
\]

Since

\[
Var(X)=E[X^2]-(E[X])^2
\]

we obtain

\[
Cov(X,Y)=-2Var(X)
\]

\subsection{Pearson’s Correlation Coefficient}

Pearson’s correlation coefficient measures the effect of change in one variable when the other variable changes.

The formula is

\[
\rho_{XY}=\frac{Cov(X,Y)}{\sigma_X\sigma_Y}
\]

Properties:

\begin{itemize}
\item Measures strength and direction of linear relationship
\item Normalized version of covariance
\item Unit-free
\end{itemize}

Range:

\[
-1 \le \rho_{XY} \le 1
\]

Interpretation:

\begin{itemize}
\item $\rho_{XY}>0$ $\rightarrow$ positive linear trend
\item $\rho_{XY}<0$ $\rightarrow$ negative linear trend
\item $\rho_{XY}=0$ $\rightarrow$ no linear correlation
\end{itemize}

\section{Jointly Gaussian Random Variables}

\textbf{Definition}

A pair $(X,Y)$ is jointly Gaussian if its joint density has the bivariate Gaussian form and the marginals are Gaussian.

Applications:

\begin{itemize}
\item signal processing
\item communication systems
\item modeling continuous noisy measurements
\end{itemize}

\textbf{Joint Gaussian PDF}

For two variables $X$ and $Y$:

\[
f_{XY}(x,y)=\frac{1}{2\pi\sigma_X\sigma_Y\sqrt{1-\rho^2}}
\exp\left(
-\frac{1}{2(1-\rho^2)}
\left[
\left(\frac{x-\mu_X}{\sigma_X}\right)^2
+
\left(\frac{y-\mu_Y}{\sigma_Y}\right)^2
-
2\rho
\left(\frac{x-\mu_X}{\sigma_X}\right)
\left(\frac{y-\mu_Y}{\sigma_Y}\right)
\right]
\right)
\]

Where

\[
\mu_X, \mu_Y = \text{means}
\]

\[
\sigma_X, \sigma_Y = \text{standard deviations}
\]

\[
\rho = \text{Pearson correlation coefficient}
\]

Example variables:

\begin{itemize}
\item $X$: Atmospheric CO$_2$ concentration (ppm)
\item $Y$: Global surface temperature anomaly ($^\circ$C)
\end{itemize}

If

\[
\rho>0
\]

then high CO$_2$ levels tend to occur with higher temperatures.

\end{document}